% Experimental Design section (Stage II). Needs amsmath, booktabs, graphicx (already loaded in report.tex).
% Generated tables: run  python3 factorial/make_tex.py --meta factorial/plan_gce/design_meta.json \
%     --results factorial/results_gce.csv --outdir report
% Generated results: python3 factorial/analyze.py ... --outdir report ; python3 factorial/plots.py ... --outdir report
% Everything marked TODO must be filled in or deleted before submission.

\section{Experimental Design}
\label{sec:design}

\subsection{System and workload parameters}

Table~\ref{tab:params} lists the parameters that can influence the behaviour of the system under test. The
factors of the experiment are chosen among them; the others are held fixed and reported so that the experiment
can be repeated.

\begin{table}[h]
  \centering
  \caption{System and workload parameters. The six parameters in bold are the factors of the experiment.}
  \label{tab:params}
  \small
  \begin{tabular}{@{}lll@{}}
    \toprule
    Parameter & Values & Role \\
    \midrule
    \textbf{Ensemble size} & 3, 5 (1 in Stage I) & factor \\
    \textbf{forceSync} (fsync of the transaction log) & yes, no & factor \\
    \textbf{Log device} (disk type of \texttt{dataLogDir}) & pd-balanced, pd-ssd & factor \\
    \textbf{Concurrent clients} & 10, 160 & factor \\
    \textbf{Write ratio} & 0.1, 0.5 & factor \\
    \textbf{Payload size} & 128\,B, 4\,KiB & factor \\
    Server JVM heap, \texttt{tickTime}, \texttt{maxClientCnxns} & 1\,GB, 2000\,ms, unlimited & fixed \\
    Machine types and zone & n2-standard-2 servers, n2-standard-8 client, one zone & fixed \\
    Znode pool, key distribution & 200 znodes, uniform & fixed \\
    Client model & closed loop, one session per client, synchronous & fixed \\
    Run length & 5\,s warm-up + 20\,s measured, fresh pool and ensemble per block & fixed \\
    \bottomrule
  \end{tabular}
\end{table}

\subsection{Metrics}

The primary metric is \emph{throughput} (successful operations per second), because the Universal Scalability
Law of Section~\ref{sec:usl} is a model of throughput. The secondary metric is the \emph{99th percentile write
latency}: coordination services are used on the critical path of other systems, so the tail of the write path
matters as much as its mean rate. Reads are served locally by any server, so read latency is reported but is
not a response of the design.

\subsection{Factors and levels}

Each factor is tied to a hypothesis from Stage I about where the cost of the write pipeline sits
(Table~\ref{tab:factors}). Levels are placed at the two ends of the range that is meaningful in practice.
The low level of the client count (10) lies below the saturation knee measured in Stage I and the high level
(160) lies well beyond it; the write ratio 0.1 is close to the production ratio of about 10:1
cited in~\cite{mit_6824_zookeeper}, and 0.5 stresses the write pipeline.
Disabling \texttt{forceSync} is not a production setting (it risks losing acknowledged writes on a crash); it
is used here only as a controlled \emph{probe} that removes the log-sync stage from the write path.

\begin{table}[h]
  \centering
  \caption{Factors, levels and justification. ``Hard'' factors need a reconfiguration of the ensemble.}
  \label{tab:factors}
  \footnotesize
  \IfFileExists{design_factors.tex}{\input{design_factors.tex}}{\textbf{[TODO: run make\_tex.py]}}
\end{table}

\subsection{Design}

A full $2^6$ design needs 64 runs per replicate. We use a half fraction, $2^{6-1}$, with the sixth factor
generated as $F=ABCDE$, i.e.\ the defining relation $I=ABCDEF$. The design has resolution VI: every main effect
is aliased only with a five-factor interaction, and every two-factor interaction only with a four-factor
interaction; three-factor interactions are aliased in pairs and are not separated. Under the usual
sparsity-of-effects assumption that interactions of order four and higher are negligible, all six main effects
and all fifteen two-factor interactions are estimated cleanly from 32 runs per replicate
(Table~\ref{tab:alias}). With two replicates the experiment has 64 runs and $32$ degrees of freedom for the
experimental error, which gives confidence intervals on every effect. The full sign table is in
Appendix~\ref{app:sign}.

\begin{table}[h]
  \centering
  \caption{Alias structure ($I=ABCDEF$, resolution VI): main effects and two-factor interactions.}
  \label{tab:alias}
  \footnotesize
  \IfFileExists{design_alias.tex}{\input{design_alias.tex}}{\textbf{[TODO: run make\_tex.py]}}
\end{table}

\paragraph{Restricted randomization.} Ensemble size, \texttt{forceSync} and log device require restarting and
reconfiguring the servers (data and logs are wiped). Runs that share the levels of these three factors form a
\emph{block} (eight blocks of four runs per replicate), so the ensemble is rebuilt eight times per replicate and
not for every run. The order of the blocks is randomized independently in each replicate, and the order of the
runs inside a block is randomized as well. Because every replicate rebuilds the servers from scratch, the
restart-to-restart variability enters the pure-error estimate. This is not a full split-plot analysis: a proper
one would test the three hard-to-change factors against a separate whole-plot error with fewer degrees of
freedom, so the confidence intervals for those three factors may be somewhat optimistic. We therefore interpret
them by their share of the variation, and consider only large effects conclusive.

\paragraph{Protocol.} Servers and load generator run on separate virtual machines. For every run the client
records its own CPU use (a client above about 70\% would make the client, and not ZooKeeper, the bottleneck)
and the runner aborts if the client prints more than a few lines of output (per-operation logging throttles the
client; this was the cause of the artificially low throughput in our first sweep).

\subsection{Analysis method}

Let $s_{ij}$ be the $\pm1$ entry of the sign table for run $i$ and effect $j$, $\bar y_i$ the mean of the $r$
replicates of run $i$ and $n=32$ the number of runs. The effects, the allocation of variation and the
confidence intervals are
\begin{equation}
  q_j=\frac1n\sum_{i} s_{ij}\,\bar y_i,\qquad
  SS_j=n\,r\,q_j^2,\qquad
  SSE=\sum_{i,k}(y_{ik}-\bar y_i)^2,\qquad
  SST=SSE+\sum_j SS_j,
\end{equation}
\begin{equation}
  s_e^2=\frac{SSE}{n(r-1)},\qquad
  q_j \pm t_{[0.975;\,n(r-1)]}\,\frac{s_e}{\sqrt{n\,r}} .
\end{equation}
The percentage of the total variation explained by effect $j$ is $100\,SS_j/SST$. With 31 effects tested at the
95\% level, about two false positives are expected, so effects are ranked by their share of the variation and
not only by significance. We check the residuals against the predicted value and on a normal quantile plot; if
the ratio between the largest and the smallest response exceeds 10, the analysis is repeated on $\log_{10}$ of
the response (a multiplicative model).

\subsection{Hypotheses}

We fix the expected outcomes before the runs, so that the experiment can falsify them.
\begin{itemize}
  \item \textbf{H1 (log-sync stage).} Disabling \texttt{forceSync} (D) increases throughput and lowers the
    write p99 latency; the eBPF fsync probe shows (almost) no fsync calls when D is at its high level.
    If D has no effect, the contention coefficient $\alpha$ of Section~\ref{sec:usl} does not come from fsync.
  \item \textbf{H2 (writes).} The benefit of disabling fsync grows with the write ratio (interaction B$\times$D),
    since reads never touch the log.
  \item \textbf{H3 (disk).} A faster log disk (F) helps only while fsync is on (interaction D$\times$F).
  \item \textbf{H4 (quorum).} A larger ensemble (A) lowers throughput and raises the write p99 latency, and
    this cost is concentrated in writes (interaction A$\times$B).
  \item \textbf{H5 (load).} The number of clients (C) has the largest main effect on throughput, because the
    low level is below and the high level above the saturation point.
\end{itemize}

\subsection{Results of the factorial experiment}
\label{sec:factorial-results}

\textbf{[TODO after the GCE runs: replace this paragraph.]} Table~\ref{tab:effects} gives the effects on
throughput and Fig.~\ref{fig:pareto} the allocation of variation. State, in this order: (i) the percentage of
the variation explained by the largest effects and by the experimental error; (ii) for each of H1--H5 whether it
was confirmed, with the effect size and its confidence interval; (iii) any interaction that was not predicted;
(iv) what the residual plots show and whether the log transform was needed; (v) what the eBPF probes showed for
the configurations at the two levels of D and F (Section~\ref{sec:ebpf}).

\begin{table}[h]
  \centering
  \caption{Largest effects on throughput (ops/s): $q$ is the effect of moving a factor from its low to its high
    level, halved, with 95\% confidence intervals.}
  \label{tab:effects}
  \footnotesize
  \IfFileExists{effects_throughput_ops_sec.tex}{\input{effects_throughput_ops_sec.tex}}{\textbf{[TODO: run analyze.py --outdir report]}}
\end{table}

\begin{figure}[h]
  \centering
  \IfFileExists{pareto_throughput_ops_sec.png}{\includegraphics[width=0.48\textwidth]{pareto_throughput_ops_sec.png}}{\textbf{[TODO: pareto plot]}}
  \hfill
  \IfFileExists{interactions_throughput_ops_sec.png}{\includegraphics[width=0.48\textwidth]{interactions_throughput_ops_sec.png}}{}
  \caption{Left: share of the total variation of throughput explained by each effect (grey: not significant).
    Right: interaction plots of the strongest two-factor interactions.}
  \label{fig:pareto}
\end{figure}

% ---- put this in the appendix of report.tex ----
% \appendix
% \section{Sign table}\label{app:sign}
% \begin{table}[h]\centering\scriptsize
%   \caption{Sign table of the $2^{6-1}$ design with the mean throughput of the two replicates.}
%   \IfFileExists{design_signtable.tex}{\input{design_signtable.tex}}{}
% \end{table}
